\documentclass[11pt,a4paper]{article}

\usepackage[T1]{fontenc}
\usepackage[utf8]{inputenc}
\usepackage[slovene]{babel}
\usepackage{lmodern}
\usepackage{amsmath,amssymb}
\usepackage{geometry}
\usepackage{microtype}
\usepackage{enumitem}
\usepackage{fancyhdr}
\usepackage[hidelinks]{hyperref}

\geometry{left=2.35cm,right=2.35cm,top=2.15cm,bottom=2.15cm}
\setlength{\parindent}{0pt}
\setlength{\parskip}{5pt}
\setlist[enumerate]{leftmargin=*,itemsep=3pt,topsep=3pt}
\setlist[itemize]{leftmargin=*,itemsep=2pt,topsep=2pt}
\pagestyle{fancy}
\fancyhf{}
\fancyfoot[C]{\thepage}
\renewcommand{\headrulewidth}{0pt}
\setcounter{secnumdepth}{0}

\newtheorem{proposition}{Domneva}

\begin{document}
\begin{center}
    {\LARGE\bfseries Projekt 8: Inducirani 2-regularni podgrafi kubičnih grafov}\\[5pt]
    {\large Predstavitev problema in načrt projektnega dela}
\end{center}
\vspace{3pt}

\section{Namen in matematično ozadje}

Cilj projektne naloge je z računalniškimi metodami preučiti \emph{inducirane 2-regularne podgrafe kubičnih grafov}. Kubični graf
je graf, v katerem ima vsako vozlišče natanko tri sosede. Podgraf, induciran z množico vozlišč
$S$, označimo z $G[S]$; vsebuje vsa vozlišča iz $S$ in vse povezave prvotnega grafa med
njimi. Podgraf je 2-regularen, kadar ima vsako njegovo vozlišče natanko dva soseda znotraj tega podgrafa.
Tak podgraf je zato unija med seboj disjunktnih ciklov.

Osrednja količina, ki jo bomo računali, je
\begin{equation}
    c_{\mathrm{ind}}(G)
    =\max\bigl\{|S|:\ S\subseteq V(G),\ G[S]\text{ je 2-regularen}\bigr\}.
\end{equation}
Pove nam, koliko vozlišč lahko največ izberemo, da bodo vsa izbrana vozlišča sestavljala inducirane cikle.
Za kubični graf z $n=|V(G)|$ vozlišči nas bo posebej zanimala razlika
\[
    \Delta(G)=c_{\mathrm{ind}}(G)-\frac{n}{2}.
\]
Ker je vsota stopenj $3n=2|E(G)|$, imajo kubični grafi sodo število vozlišč. Iskali bomo kubične grafe,
pri katerih je $\Delta(G)$ čim manjša, in z računalniškimi poskusi preverjali domnevo iz projektnih
navodil: 

\begin{proposition}
    \label{domneva}
    Če je $G$ kubični graf reda $n$, potem $c_{\mathrm{ind}}(G) \geq \frac{n}{2}$.
\end{proposition}

Pomembno je poudariti, da računanje na končnem številu primerov samo po sebi ne pomeni
splošnega dokaza domneve, lahko pa odkrije protiprimer ali ponudi koristne ugotovitve.

\section{Metoda: celoštevilsko linearno programiranje}

Najprej bomo razvili program za natančen izračun $c_{\mathrm{ind}}(G)$. Za vsako vozlišče $v$ uvedemo binarno spremenljivko
$x_v\in\{0,1\}$, ki pove, ali smo vozlišče izbrali v množico $S$. Nato rešimo optimizacijski problem
\begin{align*}
    \text{maksimiziraj}\quad & \sum_{v\in V(G)}x_v,\\[-2pt]
    \text{pri pogojih}\quad & \sum_{u\in N(v)}x_u\geq 2x_v
    &&\text{za vsak }v\in V(G),\\[-2pt]
    & \sum_{u\in N(v)}x_u\leq 3-x_v
    &&\text{za vsak }v\in V(G).
\end{align*}
Tu $N(v)$ označuje množico sosedov vozlišča $v$. Če velja $x_v=1$, pogoja zahtevata, da ima $v$
natanko dva izbrana soseda; s tem je $G[S]$ 2-regularen. Ciljna funkcija pa poskrbi, da je
izbrana množica čim večja.

Model bomo implementirali v Pythonu, na primer s knjižnicama \texttt{NetworkX} za delo z grafi in
\texttt{PuLP} oziroma drugim reševalnikom za celoštevilsko linearno programiranje. Program bo vrnil vrednost $c_{\mathrm{ind}}(G)$ in pripadajočo
optimalno množico $S$. Pravilnost bomo najprej preverili na manjših znanih grafih: $K_4$, $K_{3,3}$, prizmi, Petersenovem
in Franklinovem grafu. Pri vsakem rezultatu bomo dodatno preverili, da imajo vsa vozlišča podgrafa $G[S]$
stopnjo dve.

\section{Računalniško raziskovanje posebnih družin grafov}

Ker navodila navajajo, da so bili kubični grafi z največ 22 vozlišči že pregledani, bomo pozornost
usmerili na večje grafe. Preizkusili bomo več posebej zanimivih družin: kubične grafe z velikim obsegom
najkrajšega cikla (\emph{girth}), snarke, posplošene Petersenove grafe $P(n,k)$ in ravninske kubične grafe.

Za grafe manjših redov bomo uporabili generator \texttt{geng},
za večje grafe in grafe z večjim obsegom najkrajšega cikla pa tudi javno dostopne zbirke,
na primer \emph{House of Graphs}. Za vsak pregledani graf bomo izračunali $c_{\mathrm{ind}}(G)$ in $\Delta(G)$ ter
zabeležili primere z najmanjšimi vrednostmi. Rezultate bomo primerjali med družinami grafov in ugotavljali, katere strukturne
lastnosti so skupne grafom s posebej majhnimi vrednostmi $\Delta(G)$.

\section{Iskanje zahtevnejših primerov}

Za večje rede, okvirno $24\leq n\leq 60$, bomo uporabili tudi hevristično iskanje s \emph{simuliranim ohlajanjem}
(angl. \emph{simulated annealing}). Začeli bomo s povezanim enostavnim kubičnim grafom, nato pa bomo s preklopi
povezav ustvarjali nove kubične grafe. Pri tem bomo pazili, da graf ostane enostaven, povezan in
3-regularen.

Cilj iskanja bo zmanjševanje vrednosti $\Delta(G)$. Ker ta pri fiksnem $n$ zavzame le malo celoštevilskih
vrednosti, bomo dodali še pomožni kriterij, na primer število kratkih induciranih ciklov, da bo
iskanje lažje razlikovalo med podobno dobrimi kandidati. Za vsak obravnavani red grafa bomo poskusili poiskati
najmanjšo doseženo vrednost $\Delta(G)$, shranili ustrezne grafe in jih primerjali s trilobiti ter Franklinovim grafom.

\section{Način primerjave in preverjanje zanesljivosti}

Pri rezultatih bomo dosledno ločevali med \emph{natančnim izračunom} optimuma za posamezen graf in \emph{hevristično najdenim}
grafom z majhno vrednostjo $\Delta(G)$. Simulirano ohlajanje ne zagotavlja, da smo našli najboljši graf med
vsemi kubičnimi grafi danega reda. Zato bomo poskuse z različnimi začetnimi grafi in naključnimi semeni
večkrat ponovili, zanimive kandidate pa vedno dodatno ovrednotili z natančnim celoštevilskim modelom.

Za primerljivost poskusov bomo shranjevali red grafa, njegovo družino, vrednosti $c_{\mathrm{ind}}(G)$ in $\Delta(G)$, izbrano množico
vozlišč $S$ ter po možnosti tudi čas računanja. Posebej pomembne grafe bomo shranili v kompaktni
obliki, na primer v zapisu \texttt{graph6}, kar bo omogočalo ponovno preverjanje in prikaz najdenih primerov.

\section{Pričakovani rezultat in poročilo}

Končni rezultat projekta bo delujoč program za izračun $c_{\mathrm{ind}}(G)$, skupaj z računalniškimi poskusi na izbranih
družinah grafov in z rezultati hevrističnega iskanja. Pripraviti želimo preglednico najmanjših najdenih vrednosti $\Delta(G)$ po
številu vozlišč, opise oziroma zapise najbolj zanimivih grafov ter primerjavo s poznanimi primeri. Če bomo
našli protiprimer domnevi \ref{domneva}, ga bomo posebej izpostavili in preverili z natančnim optimizacijskim postopkom. 

% Po želji bomo za izbrane grafe dodatno izračunali še $(1,2)$-dominacijsko število $\gamma_{1,2}(G)$ in ga primerjali s $c_{\mathrm{ind}}(G)$.

\end{document}